% ============================================================
%  Příprava na maturitu – Mechatronika
%  Kompiluj: pdflatex maturita_mechatronika.tex
%  Kódování: UTF-8
% ============================================================

\documentclass[a4paper,10pt]{article}

% --- Balíčky ---
\usepackage[T1]{fontenc}
\usepackage[utf8]{inputenc}
\usepackage[czech]{babel}
\usepackage[left=2.4cm, right=1.8cm, top=1.8cm, bottom=1.8cm]{geometry}
\usepackage{lmodern}
\usepackage{microtype}
\usepackage{parskip}
\usepackage{titlesec}
\usepackage{enumitem}
\usepackage{xcolor}
\usepackage{hyperref}
\usepackage{tabularx}
\usepackage{booktabs}
\usepackage{array}
\usepackage{tikz}
\usepackage{eso-pic}
\usepackage{mdframed}
\usepackage{amsmath}
\usepackage{amssymb}
\usepackage{pgfplots}
\pgfplotsset{compat=1.18}
\usetikzlibrary{arrows.meta, positioning, shapes.geometric, calc}

% --- Barvy ---
\definecolor{accent}{HTML}{03254E}
\definecolor{stripeblue}{HTML}{568EA3}
\definecolor{medgray}{HTML}{888888}
\definecolor{lightblue}{HTML}{EAF2F8}
\definecolor{lightyellow}{HTML}{FDFBE4}
\definecolor{lightgreen}{HTML}{EAF7EA}

% --- Modrý pruh vlevo ---
\AddToShipoutPictureBG{%
  \begin{tikzpicture}[remember picture, overlay]
    \fill[stripeblue]
      (current page.north west)
      rectangle
      ([xshift=10mm] current page.south west);
    \fill[accent!60]
      ([xshift=10mm] current page.north west)
      rectangle
      ([xshift=12mm] current page.south west);
  \end{tikzpicture}%
}

\hypersetup{colorlinks=true, urlcolor=accent, linkcolor=accent}

\titleformat{\section}
  {\large\bfseries\color{accent}}
  {\thesection.}
  {0.5em}
  {}
  [\color{accent}\titlerule]

\titlespacing{\section}{0pt}{16pt}{6pt}

\newmdenv[
  backgroundcolor=lightblue,
  linecolor=stripeblue,
  linewidth=1pt,
  roundcorner=4pt,
  innerleftmargin=8pt,
  innerrightmargin=8pt,
  innertopmargin=6pt,
  innerbottommargin=6pt,
  skipabove=4pt,
  skipbelow=4pt
]{pojmybox}

\newmdenv[
  backgroundcolor=lightyellow,
  linecolor=accent!40,
  linewidth=1pt,
  roundcorner=4pt,
  innerleftmargin=8pt,
  innerrightmargin=8pt,
  innertopmargin=6pt,
  innerbottommargin=6pt,
  skipabove=4pt,
  skipbelow=8pt
]{vysvetlenibox}

\newmdenv[
  backgroundcolor=lightgreen,
  linecolor=accent!30,
  linewidth=1pt,
  roundcorner=4pt,
  innerleftmargin=8pt,
  innerrightmargin=8pt,
  innertopmargin=6pt,
  innerbottommargin=6pt,
  skipabove=4pt,
  skipbelow=8pt
]{vzorcebox}

\pagestyle{plain}

% ============================================================
\begin{document}

% --- Záhlaví ---
\begin{minipage}[t]{0.75\textwidth}
    {\Huge\bfseries\color{accent} Příprava na maturitu}\\[4pt]
    {\large\color{stripeblue} Mechatronika}\\[2pt]
    \textcolor{medgray}{\small Suchánek Lukáš \quad SPŠ Prosek \quad 2026}
\end{minipage}
\hfill
\begin{minipage}[t]{0.22\textwidth}
    \raggedleft
    \begin{tikzpicture}
        \draw[color=stripeblue, rounded corners=4pt, line width=1pt]
            (0,0) rectangle (3,1.2)
            node[midway, align=center, color=accent, font=\small\bfseries]
            {Otázek: 22};
    \end{tikzpicture}
\end{minipage}

\vspace{6pt}
\noindent\color{accent}\rule{\linewidth}{1.5pt}
\vspace{2pt}

% ============================================================
\section{Průmyslové komunikace -- komunikační sběrnice, SCADA systémy}

\begin{pojmybox}
\textbf{Klíčové pojmy:}
Profibus, Profinet, EtherCAT, Modbus, CAN, IO-Link,
OPC-UA, SCADA, HMI, PLC, topologie sítě,
master/slave, deterministická komunikace, cyklický čas
\end{pojmybox}

\begin{vysvetlenibox}
\textbf{Srovnání průmyslových sběrnic}

\begin{center}
\begin{tabularx}{0.95\linewidth}{@{} l l l X @{}}
    \toprule
    \textbf{Sběrnice} & \textbf{Přenosová rychlost}
                      & \textbf{Topologie} & \textbf{Použití} \\
    \midrule
    Modbus RTU  & 115 kbps   & linka     & Jednoduchá SCADA, měřicí přístroje \\
    Modbus TCP  & 100 Mbps   & hvězda    & Ethernet varianta Modbusu \\
    Profibus DP & 12 Mbps    & linka     & Siemens průmyslová automatizace \\
    Profinet    & 100 Mbps   & hvězda    & Nástupce Profibus přes Ethernet \\
    EtherCAT    & 100 Mbps   & prsten    & Ultrarychlé řízení pohonů (\textless{}100\,µs) \\
    CAN         & 1 Mbps     & linka     & Automobily, průmysl (kratší sítě) \\
    IO-Link     & 230 kbps   & bod-bod   & Inteligentní senzory/aktuátory \\
    \bottomrule
\end{tabularx}
\end{center}

\textbf{Model OSI a průmyslové protokoly}\\
Průmyslové sběrnice obvykle implementují pouze vrstvy 1 (fyzická),
2 (linková) a 7 (aplikační) -- mezivrstvy jsou vynechány pro minimalizaci
latence.

\textbf{SCADA systém -- architektura}

\begin{center}
\begin{tikzpicture}[scale=0.85,
  lvl/.style={rectangle, draw=accent, fill=lightblue, rounded corners=3pt,
              minimum width=7cm, minimum height=0.65cm,
              text centered, font=\small}]
  \node[lvl, fill=lightyellow] (l4) at (0,3)   {Úroveň 4: ERP (podnikové řízení)};
  \node[lvl]                   (l3) at (0,2.2) {Úroveň 3: MES (řízení výroby)};
  \node[lvl, fill=lightgreen]  (l2) at (0,1.4) {Úroveň 2: SCADA / HMI (vizualizace)};
  \node[lvl]                   (l1) at (0,0.6) {Úroveň 1: PLC / DCS (lokální řízení)};
  \node[lvl, fill=lightyellow] (l0) at (0,-0.2){Úroveň 0: Senzory, aktuátory, pohony};
  \foreach \a/\b in {l4/l3, l3/l2, l2/l1, l1/l0}
    \draw[>-, accent] (\a) -- (\b);
\end{tikzpicture}
\end{center}

\textbf{OPC-UA} -- moderní, platformově nezávislý komunikační standard
pro vertikální integraci (od senzoru až po ERP), šifrovaný, s datovými
modely; základ Industry~4.0 komunikace.
\end{vysvetlenibox}

% ============================================================

% ============================================================
\end{document}
